\documentclass[11pt]{article}

\usepackage[margin=1in]{geometry}
\usepackage{amsmath}
\usepackage{amssymb}

\setlength{\parindent}{0pt}
\setlength{\parskip}{0.6em}

\begin{document}

\begin{center}
    {\large\textbf{STAT GU4204 --- Statistical Inference}}\\[2pt]
    {\normalsize Homework 3 --- Full Solutions}\\[2pt]
    {\small Jonah Aden \quad $\cdot$ \quad Columbia University, Spring 2026}
\end{center}

\vspace{1em}

%%%%%%%%%%%%%%%%%%%%%%%%%%%%%%%%%%%%%%%%%%%%%%%%%%%%%%%%%%%%%%%%%%%%%%%%%%%%%%%
\section*{Problem 8.5.2}
\textbf{Setup.} The eight observations $X_1,\dots,X_8 \stackrel{\text{iid}}{\sim} N(\mu,\sigma^2)$ with $\mu$ and $\sigma^2$ both unknown. The standard pivot for $\mu$ in this situation is
\[
T = \frac{\bar{X}_n - \mu}{S'/\sqrt{n}} \sim t_{n-1},
\qquad
S'^2 := \frac{1}{n-1}\sum_{i=1}^n (X_i - \bar{X}_n)^2,
\]
which gives the symmetric two-sided CI
\[
\bar{X}_n \;\pm\; t_{(1+\gamma)/2,\,n-1}\,\frac{S'}{\sqrt{n}}.
\]
Because the $t$-density is symmetric about $0$ and unimodal, the symmetric (equal-tailed) interval is also the \emph{shortest} confidence interval at any given confidence level.

\textbf{Sample statistics.} With the data $3.1, 3.5, 2.6, 3.4, 3.8, 3.0, 2.9, 2.2$:
\[
\sum_i x_i = 24.5
\quad\Longrightarrow\quad
\bar{x} = \frac{24.5}{8} = 3.0625.
\]

The squared deviations are
\begin{align*}
(0.0375)^2 &= 0.00140625 \\
(0.4375)^2 &= 0.19140625 \\
(-0.4625)^2 &= 0.21390625 \\
(0.3375)^2 &= 0.11390625 \\
(0.7375)^2 &= 0.54390625 \\
(-0.0625)^2 &= 0.00390625 \\
(-0.1625)^2 &= 0.02640625 \\
(-0.8625)^2 &= 0.74390625
\end{align*}
giving $\sum_i (x_i - \bar{x})^2 = 1.83875$. Therefore
\[
s'^2 = \frac{1.83875}{7} = 0.26268,
\qquad
s' = 0.5125,
\qquad
\frac{s'}{\sqrt{8}} = \frac{0.5125}{2.8284} = 0.1812.
\]

\textbf{Critical values from the $t_7$ table.}
\[
t_{0.95,\,7} = 1.895, \qquad
t_{0.975,\,7} = 2.365, \qquad
t_{0.995,\,7} = 3.499.
\]

\textbf{Confidence intervals.}
\begin{align*}
\text{(a) } 90\%: \quad & 3.0625 \pm (1.895)(0.1812) = 3.0625 \pm 0.3434 = \boxed{(2.719,\ 3.406)}. \\[4pt]
\text{(b) } 95\%: \quad & 3.0625 \pm (2.365)(0.1812) = 3.0625 \pm 0.4285 = \boxed{(2.634,\ 3.491)}. \\[4pt]
\text{(c) } 99\%: \quad & 3.0625 \pm (3.499)(0.1812) = 3.0625 \pm 0.6340 = \boxed{(2.428,\ 3.697)}.
\end{align*}

%%%%%%%%%%%%%%%%%%%%%%%%%%%%%%%%%%%%%%%%%%%%%%%%%%%%%%%%%%%%%%%%%%%%%%%%%%%%%%%
\section*{Problem 8.5.4}

\textbf{Setup.} Because $\sigma^2$ is known, we use the standard normal pivot $Z = (\bar{X}_n - \mu)/(\sigma/\sqrt{n}) \sim N(0,1)$, giving the symmetric $95\%$ CI
\[
\bar{X}_n \pm z_{0.975}\,\frac{\sigma}{\sqrt{n}},
\qquad
\text{length } L = 2\,z_{0.975}\,\frac{\sigma}{\sqrt{n}}.
\]
With $z_{0.975} = 1.96$, the requirement $L \le 0.01\sigma$ becomes
\[
\frac{2(1.96)\sigma}{\sqrt{n}} \le 0.01\sigma
\quad\Longleftrightarrow\quad
\sqrt{n} \ge \frac{2(1.96)}{0.01} = 392
\quad\Longleftrightarrow\quad
n \ge 392^2 = 153{,}664.
\]

\textbf{Answer.} The smallest sample size that produces a $95\%$ CI of length $\le 0.01\sigma$ is
\[
\boxed{n \ge 153{,}664.}
\]

%%%%%%%%%%%%%%%%%%%%%%%%%%%%%%%%%%%%%%%%%%%%%%%%%%%%%%%%%%%%%%%%%%%%%%%%%%%%%%%
\section*{Problem 8.7.11}

Let $\sigma_B^2 = \sigma^2$, so $\sigma_A^2 = 4\sigma^2$. Then
\[
\operatorname{Var}(\bar{X}_m) = \frac{4\sigma^2}{m},
\qquad
\operatorname{Var}(\bar{Y}_n) = \frac{\sigma^2}{n},
\qquad
\bar{X}_m \perp \bar{Y}_n.
\]

\textbf{(a) Unbiasedness.} For any constants $\alpha, m, n$,
\[
E[\hat{\theta}] = \alpha\,E[\bar{X}_m] + (1-\alpha)\,E[\bar{Y}_n]
= \alpha\theta + (1-\alpha)\theta = \theta.
\]
Therefore $\hat{\theta}$ is unbiased for $\theta$ \emph{for every} $\alpha \in \mathbb{R}$ and every $m, n \ge 1$.

\textbf{(b) Minimum-variance $\alpha$.} By independence,
\[
\operatorname{Var}(\hat{\theta})
= \alpha^2 \cdot \frac{4\sigma^2}{m} + (1-\alpha)^2 \cdot \frac{\sigma^2}{n}
= \sigma^2 \left[\frac{4\alpha^2}{m} + \frac{(1-\alpha)^2}{n}\right].
\]
Differentiating with respect to $\alpha$ and setting to zero:
\[
\frac{d}{d\alpha}\left[\frac{4\alpha^2}{m} + \frac{(1-\alpha)^2}{n}\right]
= \frac{8\alpha}{m} - \frac{2(1-\alpha)}{n} = 0.
\]
Multiplying through by $mn$ and rearranging:
\[
8\alpha n = 2m(1-\alpha)
\;\Longrightarrow\;
8\alpha n + 2\alpha m = 2m
\;\Longrightarrow\;
\alpha(8n + 2m) = 2m
\;\Longrightarrow\;
\boxed{\alpha^* = \frac{m}{m + 4n}.}
\]
The second derivative is $8/m + 2/n > 0$, confirming a global minimum.

%%%%%%%%%%%%%%%%%%%%%%%%%%%%%%%%%%%%%%%%%%%%%%%%%%%%%%%%%%%%%%%%%%%%%%%%%%%%%%%
\section*{Problem 8.8.2}

DeGroot \S5.5 defines the geometric distribution as the number of failures before the first success, with PMF
\[
f(x \mid p) = p(1-p)^x, \qquad x = 0, 1, 2, \dots,
\quad E[X] = \frac{1-p}{p}.
\]

\textbf{Score and curvature.}
\[
\log f(x\mid p) = \log p + x \log(1-p),
\qquad
\frac{\partial}{\partial p}\log f = \frac{1}{p} - \frac{x}{1-p}.
\]
\[
\frac{\partial^2}{\partial p^2} \log f = -\frac{1}{p^2} - \frac{x}{(1-p)^2}.
\]

\textbf{Fisher information.} Using $E[X] = (1-p)/p$,
\[
I(p) = -E\!\left[\frac{\partial^2}{\partial p^2}\log f\right]
= \frac{1}{p^2} + \frac{E[X]}{(1-p)^2}
= \frac{1}{p^2} + \frac{(1-p)/p}{(1-p)^2}
= \frac{1}{p^2} + \frac{1}{p(1-p)}.
\]
Combining over the common denominator $p^2(1-p)$:
\[
I(p) = \frac{(1-p) + p}{p^2(1-p)} = \boxed{\frac{1}{p^2(1-p)}.}
\]

%%%%%%%%%%%%%%%%%%%%%%%%%%%%%%%%%%%%%%%%%%%%%%%%%%%%%%%%%%%%%%%%%%%%%%%%%%%%%%%
\section*{Problem 8.8.7}

For Bernoulli$(p)$, $f(x \mid p) = p^x(1-p)^{1-x}$, $x\in\{0,1\}$, with $E[X] = p$ and $\operatorname{Var}(X) = p(1-p)$.

\textbf{Fisher information per observation.}
\[
\log f = x\log p + (1-x)\log(1-p),
\qquad
\frac{\partial^2}{\partial p^2}\log f = -\frac{x}{p^2} - \frac{1-x}{(1-p)^2}.
\]
\[
I_1(p) = E\!\left[\frac{X}{p^2} + \frac{1-X}{(1-p)^2}\right]
= \frac{p}{p^2} + \frac{1-p}{(1-p)^2}
= \frac{1}{p} + \frac{1}{1-p}
= \frac{1}{p(1-p)}.
\]
For an iid sample of size $n$, the information adds: $I_n(p) = n/[p(1-p)]$.

\textbf{Cram\'er--Rao lower bound.} For any unbiased estimator $T(X_1,\dots,X_n)$ of $p$,
\[
\operatorname{Var}(T) \ge \frac{1}{I_n(p)} = \frac{p(1-p)}{n}.
\]

\textbf{Performance of $\bar{X}_n$.} Since $\bar{X}_n$ is the sample mean of iid Bernoullis with mean $p$ and variance $p(1-p)$,
\[
E[\bar{X}_n] = p \quad \text{(unbiased)},
\qquad
\operatorname{Var}(\bar{X}_n) = \frac{p(1-p)}{n}.
\]
This exactly matches the CRLB, so $\bar{X}_n$ attains the lower bound:
\[
\boxed{\bar{X}_n \text{ is an efficient estimator of } p.}
\]

%%%%%%%%%%%%%%%%%%%%%%%%%%%%%%%%%%%%%%%%%%%%%%%%%%%%%%%%%%%%%%%%%%%%%%%%%%%%%%%
\section*{Problem 8.8.17}

For $X \sim \mathrm{Bin}(n, p)$ with $n$ known,
\[
f(x \mid p) = \binom{n}{x} p^x(1-p)^{n-x},
\qquad
\log f = \log\binom{n}{x} + x\log p + (n-x)\log(1-p).
\]

\textbf{Score and curvature.}
\[
\frac{\partial}{\partial p}\log f = \frac{x}{p} - \frac{n-x}{1-p},
\qquad
\frac{\partial^2}{\partial p^2}\log f = -\frac{x}{p^2} - \frac{n-x}{(1-p)^2}.
\]

\textbf{Fisher information.} Using $E[X] = np$ and $E[n-X] = n(1-p)$,
\[
I(p) = -E\!\left[\frac{\partial^2}{\partial p^2}\log f\right]
= \frac{E[X]}{p^2} + \frac{E[n-X]}{(1-p)^2}
= \frac{np}{p^2} + \frac{n(1-p)}{(1-p)^2}
= \frac{n}{p} + \frac{n}{1-p}.
\]
Combining,
\[
I(p) = \frac{n[(1-p) + p]}{p(1-p)} = \boxed{\frac{n}{p(1-p)}.}
\]

\textbf{Remark.} This is exactly $n$ times the Bernoulli information $1/[p(1-p)]$ from Problem 8.8.7, as expected since a $\mathrm{Bin}(n,p)$ random variable has the same information as $n$ iid Bernoulli$(p)$ observations.

%%%%%%%%%%%%%%%%%%%%%%%%%%%%%%%%%%%%%%%%%%%%%%%%%%%%%%%%%%%%%%%%%%%%%%%%%%%%%%%
\section*{Problem 9.1.1(a)}

DeGroot parameterizes the exponential distribution with rate $\beta$:
\[
f(x \mid \beta) = \beta e^{-\beta x}, \qquad x \ge 0,\ \beta > 0,
\qquad
P_\beta(X \ge c) = e^{-\beta c} \text{ for } c \ge 0.
\]

The test rejects $H_0$ exactly when $X \ge 1$, so the power function is
\[
\pi(\beta) = P_\beta(X \ge 1) = \int_1^\infty \beta e^{-\beta x}\,dx
= \left[-e^{-\beta x}\right]_1^\infty = e^{-\beta}.
\]

\[
\boxed{\pi(\beta) = e^{-\beta}, \quad \beta > 0.}
\]

\textbf{Brief sanity check.} The size of the test is
\[
\sup_{\beta \ge 1} \pi(\beta) = \sup_{\beta \ge 1} e^{-\beta} = e^{-1} \approx 0.368,
\]
attained at $\beta = 1$. Power increases as $\beta$ decreases (i.e., as the alternative becomes more pronounced), which is the qualitatively correct behavior for a sensible test of $H_1 : \beta < 1$.

\end{document}
